\documentclass[11pt]{article}
\usepackage[letterpaper,margin=0.78in]{geometry}
\usepackage{amsmath,amssymb,mathtools}
\usepackage{booktabs,longtable,array,tabularx}
\usepackage{enumitem}
\usepackage{microtype}
\usepackage{xcolor}
\usepackage{hyperref}
\usepackage{xurl}
\usepackage{fancyhdr}
\usepackage{lastpage}
\usepackage{siunitx}
\usepackage{graphicx}
\usepackage{float}
\usepackage[T1]{fontenc}
\usepackage{lmodern}
\definecolor{navy}{RGB}{25,54,93}
\definecolor{slate}{RGB}{75,85,99}
\hypersetup{colorlinks=true,linkcolor=navy,urlcolor=navy,citecolor=navy,pdftitle={AWR2944 FMCW/TWTT MATLAB Simulation Blueprint}}
\setlist{nosep,leftmargin=*}
\setlist[description]{style=nextline}
\setlength{\parindent}{0pt}
\setlength{\parskip}{0.52em}
\pagestyle{fancy}
\fancyhf{}
\lhead{\small FMCW / Two-Way Time Transfer Research Corpus}
\rhead{\small AWR2944 FMCW/TWTT MATLAB Simulation Blueprint}
\cfoot{\small Page \thepage\ of \pageref{LastPage}}
\newcommand{\Saeed}{Saeed}
\newcommand{\AWR}{AWR2944}
\newcommand{\DCA}{DCA1000}
\newcommand{\TWTT}{two-way time transfer}
\newcommand{\FMCW}{FMCW}
\newcommand{\codeword}[1]{\path{#1}}
\newenvironment{keybox}{\begin{center}\begin{minipage}{0.94\textwidth}\hrule\vspace{0.45em}\textbf{Key point.} }{\vspace{0.45em}\hrule\end{minipage}\end{center}}
\begin{document}
\begin{titlepage}
\centering
\vspace*{1.2in}
{\Huge\bfseries\color{navy} AWR2944 FMCW/TWTT MATLAB Simulation Blueprint\par}
\vspace{0.3in}
{\Large A staged implementation plan from analytic truth model to two-node nonideal digital twin\par}
\vspace{0.6in}
{\large Summer 2026 - Saeed FMCW/TWTT Project\par}
\vfill
{\large Curated research synthesis and engineering reference\par}
{\large Prepared 10 August 2026\par}
\vspace{0.4in}
\begin{minipage}{0.82\textwidth}
\small
This document is a project research aid, not a claim of novelty. It distinguishes source-supported prior art, engineering inference, and proposed simulation steps. Publisher/IEEE PDFs supplied by the user are retained in the archive for private research and are not represented as redistributable.
\end{minipage}
\end{titlepage}
\tableofcontents
\clearpage

\section{Design objective}
Build one MATLAB codebase that can serve three roles without rewriting the physics: (1) a transparent ideal FMCW demonstration for Saeed/program-manager plots; (2) an ideal reciprocal two-way timing simulator that recovers propagation delay and clock offset; and (3) a progressively more realistic AWR2944-oriented precision simulator with clock, PLL, front-end, sampling, and coding impairments.

The design must prioritize verifiability over feature count. Every block needs a closed-form limiting case or a unit test. A block that cannot be independently tested should not be added to V0/V1.

\section{Why complex baseband is the correct starting representation}
A 77-GHz carrier does not need to be explicitly sampled to reproduce FMCW dechirping physics. Represent the chirp by its analytic phase, preserve exact fractional delays, and evaluate the mixer difference phase. This captures the information that reaches the IF/ADC while keeping the simulation sample rate tied to IF bandwidth rather than RF carrier frequency.

This approach also matches the philosophy of the Scheiblhofer MATLAB simulator: simulate on a time grid sufficient for IF/noise dynamics, then decimate to the final ADC rate. The first version can simplify further by generating the dechirped analytic signal exactly; later versions can insert oversampled IF filters and ADC decimation.

\section{Repository/code layout}
Recommended MATLAB layout:
\begin{verbatim}
10_SIMULATION_WORKSPACE/
  README.md
  config/
    cfg_ideal_demo.m
    cfg_awr_smoke_v1.m
  src/
    waveform/
      fmcw_phase.m
      fmcw_chirp.m
      apply_fractional_delay.m
    clocks/
      local_time_map.m
      clock_offset_model.m
    channel/
      free_space_delay.m
      multipath_channel.m
    rf/
      dechirp.m
      if_filter.m
      quantize_adc.m
    impairments/
      add_cfo.m
      add_slope_error.m
      add_phase_noise_psd.m
      add_chirp_nonlinearity.m
      add_group_delay.m
    estimation/
      estimate_fft_peak.m
      estimate_phase_slope.m
      estimate_czt.m
      solve_twtt.m
    plots/
      plot_chirp_geometry.m
      plot_beat_spectrum.m
      plot_twtt_summary.m
  tests/
    test_ideal_delay.m
    test_fractional_delay.m
    test_twtt_recovery.m
    test_cfo_up_down.m
    test_slope_mismatch.m
  demos/
    demo01_single_link_ideal.m
    demo02_twtt_ideal.m
    demo03_estimator_resolution.m
    demo04_awr_parameters.m
\end{verbatim}

\section{Configuration model}
Use one explicit struct. Avoid hidden globals and toolbox-default parameters.
\begin{verbatim}
cfg.c = 299792458;
cfg.f0 = 77e9;
cfg.B = ...;
cfg.Tc = ...;
cfg.S = cfg.B/cfg.Tc;
cfg.Fs = 10e6;
cfg.N = 256;

cfg.link.tau = 5e-9;
cfg.clock.theta = 100e-12;
cfg.clock.skew = 0;
cfg.clock.cfo = 0;

cfg.impair.awgn_snr_db = Inf;
cfg.impair.phase_noise_enable = false;
cfg.impair.nonlinearity_enable = false;
cfg.impair.group_delay_enable = false;
\end{verbatim}
Convert every TI unit (GHz, MHz/us, ksps, us) at the configuration boundary. Internal calculations should remain SI.

\section{V0: single-link analytic truth model}
\subsection{Inputs}
$f_0$, $S$, observation interval, one-way delay $\tau$, complex amplitude $\alpha$, optional initial phase.

\subsection{Waveform}
Generate phase
\[
\phi(t)=\phi_0+2\pi f_0t+\pi St^2
\]
and chirp $s(t)=e^{j\phi(t)}$. Apply the delay by evaluating $\phi(t-\tau)$, not by shifting the sample array. This keeps 10-ps-class test delays exact even when ADC sample spacing is tens or hundreds of nanoseconds.

\subsection{Dechirp}
With the selected convention,
\[
z(t)=s(t)r^*(t).
\]
The no-noise phase should be linear. Its slope must equal $2\pi S\tau$ to numerical precision.

\subsection{Estimators}
Implement two estimators immediately:
\begin{enumerate}
\item FFT peak for intuitive plotting and coarse sanity checks.
\item phase-slope least-squares estimator for transparent sub-bin frequency recovery.
\end{enumerate}
The FFT is not allowed to be the ground-truth estimator in tests.

\subsection{V0 acceptance criteria}
\begin{itemize}
\item Relative error in recovered $f_b$ below a numerical tolerance for noise-free cases.
\item Fractional delays smaller than one ADC sample recovered correctly.
\item Doubling $S$ doubles the estimated beat.
\item Switching mixer conjugation changes sign only.
\end{itemize}

\section{V0 plotting package}
\subsection{Figure A: chirp geometry}
Plot instantaneous frequency versus time for local TX and delayed RX. Annotate horizontal frequency separation $S\tau$. The plot should visually explain why a time shift becomes a frequency offset on a linear slope.

\subsection{Figure B: dechirped signal}
Show a short time trace of the complex beat phase or real IF plus a spectrum. Annotate theoretical and estimated $f_b$.

\subsection{Figure C: delay sweep}
Sweep $\tau$ over a reasonable range and plot estimated $f_b$ versus $\tau$. Overlay the theoretical line $S\tau$. This is an extremely strong program-manager plot because it shows linearity and validates the model over more than one hand-picked point.

\section{V1: ideal reciprocal two-way FMCW timing}
Create station A and B with equal rate/slope and a relative epoch offset $\theta$. Do not yet simulate a separate physical PLL. Instead define the effective relative delays:
\[
\delta_{AB}=\tau+\theta,\qquad
\delta_{BA}=\tau-\theta.
\]
Generate two dechirped observations using the same waveform functions. Estimate signed $f_{AB}$ and $f_{BA}$, then solve
\[
\hat\tau=\frac{f_{AB}+f_{BA}}{2S},\qquad
\hat\theta=\frac{f_{AB}-f_{BA}}{2S}.
\]

\subsection{Recommended demonstration numbers}
Use the AWR-like slope $S=29.982\,\mathrm{MHz}/\mu\mathrm{s}$ with $\tau=5$ ns and $\theta=100$ ps. The resulting beats are separated by only about 6 kHz, making the point that sub-bin estimation matters without requiring absurdly tiny frequencies.

\subsection{V1 acceptance criteria}
\begin{itemize}
\item $\hat\tau$ and $\hat\theta$ recover injected values to noise-free numerical tolerance.
\item Sign conventions are documented in code comments and unit tests.
\item If $\theta=0$, two beats are equal.
\item If $\tau=0$, beats are equal-and-opposite under the ideal equation.
\end{itemize}

\section{V1 program-manager figure}
Use either one composite figure or two clean figures:
\begin{itemize}
\item top: A-to-B and B-to-A beat spectra/time-frequency markers;
\item bottom/right: a small equation box reporting injected and recovered TOF and clock offset.
\end{itemize}
Avoid showing dozens of radar DSP panels. The point is the transformation ``two tiny time quantities $\rightarrow$ two easy beat frequencies $\rightarrow$ sum/difference recovery.''

\section{V1.5: estimator-resolution demonstration}
Before adding hardware noise, demonstrate why nearest-bin FFT is insufficient. For the current short ADC window, show:
\begin{enumerate}
\item raw FFT with the true beat between bins;
\item nearest-bin estimate;
\item zero-padded/interpolated display;
\item phase-slope estimate;
\item later CZT estimate.
\end{enumerate}
Plot timing error in ps for each. This turns the ``10 ps vs 39-kHz FFT bin'' concern into a concrete signal-processing story.

\section{V2: independent time and frequency offsets}
Add a carrier-frequency offset $\Delta f$ between the local and received chirps. For an up-chirp,
\[
f_+=\Delta f+S\tau_{eff}.
\]
For a down-chirp,
\[
f_-=\Delta f-S\tau_{eff}.
\]
Use both to recover $\Delta f$ and $\tau_{eff}$. This reproduces the conceptual method in the 2007 cooperative synchronization paper and gives the model its first real ``clock-rate/frequency'' dimension.

\section{V2.1: clock skew and sample-rate offset}
Represent relative epoch as $\theta(t)=\theta_0+\epsilon t$. Separately allow ADC sample-rate error $F_{s,B}=F_s(1+\epsilon_s)$. Verify their signatures differ. The first creates physical/time-map evolution; the second distorts the digital time axis. In real hardware they may be related to the same reference, but the model should not assume that until the architecture is verified.

\section{V2.2: chirp slope mismatch}
Set $S_B=S_A+\Delta S$. The dechirped signal should become chirped rather than a pure tone. Unit test this by estimating instantaneous phase slope and verifying the predicted linear drift. This is important because a two-node estimator that assumes a stationary beat can silently report biased timing when slopes differ.

\section{V3: white-noise Monte Carlo}
Add post-dechirp complex AWGN first. Sweep SNR and run at least $10^3$ Monte Carlo trials per operating point for initial plots, increasing as needed for tails. Record bias, standard deviation, RMS error, estimator failure rate, and computational time.

Suggested plots:
\begin{itemize}
\item $\sigma_{\hat\tau}$ versus SNR;
\item $\sigma_{\hat\theta}$ versus SNR;
\item estimator comparison: FFT/interpolated FFT/phase-slope/CZT;
\item precision versus observation duration/chirp count;
\item precision versus slope $S$ at fixed SNR.
\end{itemize}

\section{V4: phase noise}
Do not begin with arbitrary phase jitter samples. Implement two modes:
\begin{description}
\item[Simple validation mode.] White frequency/phase process with a known analytic behavior, used only to validate plumbing.
\item[PSD mode.] Generate colored phase noise whose one-sided/two-sided PSD convention is explicit. Use measured or documented masks later.
\end{description}

For monostatic/homodyne validation, use a common phase-noise realization for TX and LO so that the range-correlation difference $\phi(t)-\phi(t-\tau)$ emerges. For two independent stations, use separate realizations. For CFDDS-inspired experiments, process the two reciprocal observations jointly and measure how much of the noise-like term cancels.

\section{V5: systematic chirp distortion}
Implement deterministic $e_f(t)$ or $\phi_{sys}(t)$ functions:
\begin{enumerate}
\item quadratic curvature;
\item cubic curvature;
\item sinusoidal spur/ripple;
\item later, an imported measured residual trace.
\end{enumerate}
For each, show beat spectrogram/phase residual and timing bias. Validate that setting the coefficients to zero reproduces V3 exactly.

\section{V6: front-end and ADC fidelity}
Following the Scheiblhofer block philosophy, introduce:
\begin{itemize}
\item TX/RX complex gain;
\item leakage/direct coupling;
\item IF low-pass/band-pass transfer function;
\item per-path group delay;
\item ADC clipping/full scale;
\item quantization/ENOB;
\item oversampled simulation clock followed by decimation to AWR/DCA sample rate.
\end{itemize}
Keep this layer separate from oscillator impairments so each contribution to timing error can be toggled independently.

\section{V7: multipath}
Use a direct active cross-link plus one or more delayed weaker paths. Sweep relative amplitude and excess delay. Evaluate bias of frequency-only, phase-only, and combined estimators. This is a likely source of significant real-world error in parking-lot/indoor tests.

\section{V8: AWR2944-faithful configuration}
Create a config parser or manual mapping from the actual profile used by the radar package. At minimum capture:
\begin{itemize}
\item start frequency;
\item slope;
\item ADC start time;
\item sample rate and number of samples;
\item ramp end time/idle time;
\item chirps per frame and frame period;
\item TX/RX selection;
\item any known analog HPF/IF constraints relevant to low beat frequencies.
\end{itemize}

Then run the same V0/V1 tests with AWR-derived parameters. The goal is not to emulate every BSS register. It is to guarantee that the simulated waveform and acquisition window correspond to what the package actually configures.

\section{Important hardware question: very low beat frequencies}
Picosecond offsets can map to hundreds of hertz at the current slope. AWR analog HPF corners, leakage cancellation, DC offsets, and the need to avoid near-DC beat content can therefore matter. The early cooperative paper intentionally uses known frequency offsets to keep mixed signals away from DC. Before finalizing the hardware waveform, inspect AWR HPF/IF configuration and consider adding a deliberate beat-frequency offset that is algebraically removed in processing.

\section{V9: hardware-delay calibration model}
Give each station separate $g^{TX}$ and $g^{RX}$ delays plus frequency-dependent group delay. Run reciprocal solves before and after calibration. Distinguish:
\begin{itemize}
\item repeatability/precision of an unchanging setup;
\item absolute time offset accuracy after fixed-delay calibration;
\item drift with temperature, power cycle, RF channel, and configuration.
\end{itemize}
This distinction will likely become central if the goal is ``10 ps.''

\section{V10: phase-coded node identity}
Only after V1--V3 are stable, add phase coding:
\[
s_i(t)=s_{FMCW}(t)e^{j\phi_{c,i}(t)}.
\]
Start with two orthogonal or low-cross-correlation code sequences. Demonstrate that a receiver can identify A versus B at different amplitudes. Then introduce propagation delay that misaligns code chips. Compare naive decoding with an alignment/group-delay strategy inspired by Uysal/Lampel/Kumbul.

Metrics:
\begin{itemize}
\item code detection probability;
\item cross-correlation sidelobe/interference rejection;
\item timing bias after coding;
\item required ADC/IF bandwidth;
\item SNR loss relative to uncoded FMCW;
\item number of simultaneously identifiable nodes.
\end{itemize}

\section{Test philosophy}
Each physical mechanism gets three tests:
\begin{enumerate}
\item \textbf{zero test}: disabled mechanism changes nothing;
\item \textbf{analytic test}: one simple case matches a closed-form prediction;
\item \textbf{stress test}: large value produces the expected qualitative failure signature.
\end{enumerate}
Do not accept a Monte Carlo plot as evidence that a block is correct unless the zero and analytic tests pass first.

\section{Numerical precision}
Use double precision throughout V0--V5. Avoid subtracting two enormous 77-GHz phase numbers when a mathematically equivalent baseband expression can be evaluated directly; catastrophic phase wrapping can obscure small timing differences. Where absolute RF phase is needed, reduce arguments modulo $2\pi$ carefully and compare against symbolic/analytic expressions.

\section{Data products}
Every demo should save a machine-readable result struct containing config, random seed, theoretical truth, estimate, bias/error, and software version/hash. Plots should be regenerable from saved results. This is especially useful when the simulation becomes part of a paper and we need reproducibility months later.

\section{Recommended first coding sprint}
\begin{enumerate}
\item Write config + unit conversion helper.
\item Implement analytic chirp phase and delayed phase.
\item Implement ideal dechirp.
\item Implement phase-slope estimator and FFT plotter.
\item Unit-test $f_b=S\tau$ with 5 ns, 100 ps, and 10 ps cases.
\item Implement reciprocal two-way wrapper and solve $\tau,\theta$.
\item Create three clean figures.
\item Only then decide whether to add CZT before showing Saeed.
\end{enumerate}
This sprint is intentionally small enough to complete quickly while still producing a model we will keep.

\section{Definition of done for the first Saeed deliverable}
The first MATLAB deliverable is done when:
\begin{itemize}
\item equations in the code comments match the mathematical handbook;
\item all ideal unit tests pass;
\item fractional delay is not sample-quantized;
\item a clean beat-frequency plot matches theory;
\item two-way recovery returns injected TOF and clock offset;
\item the plot captions explicitly state the ideal assumptions;
\item the script runs from a clean MATLAB path with no manual variable editing.
\end{itemize}
Anything beyond this is optional for the first program-manager presentation.

\section{What comes immediately after the first plots}
The next scientifically useful task is not phase coding. It is to add independent frequency offset and up/down chirps, because that tests whether the simulator can distinguish epoch offset from oscillator mismatch in the same way the foundational cooperative work does. After that, compare beat estimators under AWGN. These two steps establish the backbone needed before phase noise and hardware calibration.

\section{Final validation against hardware}
When the second AWR2944 arrives:
\begin{enumerate}
\item characterize each board separately with the current package;
\item establish a shared or measured frequency reference strategy;
\item record direct line-of-sight cross-link data if hardware mode permits;
\item extract actual chirp/ADC parameters into the simulation config;
\item compare measured beat frequency/phase statistics to ideal simulation;
\item fit only one impairment class at a time (CFO/skew, PN, group delay, nonlinearity);
\item freeze a calibrated model before attempting coded operation.
\end{enumerate}
The simulator should become a diagnostic digital twin, not merely a figure generator.

\section{Reference anchors}
The architecture follows the modular ideas in Scheiblhofer et al. 2008; the ideal time/frequency synchronization progression follows Roehr/Vossiek/Gulden 2007; double-sided processing is benchmarked against Gottinger et al. 2019; estimator design is anchored by Scherr et al. 2015; nonlinearity/phase-noise layers follow Ayhan 2016, Herzel 2018, and Tschapek 2022; coded extensions follow Lampel/Uysal/Kumbul; AWR parameter mapping follows current TI documentation in the archive.


\clearpage
\section{Module contracts}
Each module should have a small, documented input/output contract. Suggested contracts follow.

\subsection*{\codeword{fmcw_phase(t,cfg)}}
\textbf{Input:} physical/local time vector and waveform config. \textbf{Output:} deterministic phase in radians. \textbf{Invariant:} derivative divided by $2\pi$ equals configured instantaneous frequency. \textbf{No noise inside this function.}

\subsection*{\codeword{apply_fractional_delay(t,phase,tau)}}
Prefer evaluation of the analytic phase at $t-\tau$ for V0. Later, for arbitrary sampled waveforms, use a validated fractional-delay filter. \textbf{Invariant:} a sinusoid receives phase shift $-2\pi f\tau$ without amplitude distortion over the tested band.

\subsection*{\codeword{dechirp(rx,lo)}}
Return complex product with an explicitly documented conjugation order. \textbf{Invariant:} one ideal delayed up-chirp produces a positive or negative pure tone exactly matching the selected convention.

\subsection*{\codeword{estimate_phase_slope(z,Fs)}}
Return $\hat f$, phase intercept, fit residual, and a quality flag. Reject/flag records with excessive phase residual rather than silently returning a number.

\subsection*{\codeword{solve_twtt(fAB,fBA,S)}}
Return $\hat\tau$, $\hat\theta$ plus the sign convention. This function should contain almost no DSP; it is algebra and therefore easy to test exhaustively.

\section{Detailed unit-test matrix}
\begin{longtable}{@{}p{0.05\textwidth}p{0.28\textwidth}p{0.27\textwidth}p{0.27\textwidth}@{}}
\toprule ID & Test & Injected condition & Expected result \\ \midrule
T01 & zero delay & $\tau=0$ & zero timing beat (or known offset only) \\
T02 & 5-ns delay & ideal & $f_b=S\tau$ \\
T03 & 100-ps delay & ideal & sub-bin estimator exact in noise-free case \\
T04 & 10-ps delay & ideal & no sample-quantization error \\
T05 & half-sample delay & choose $\tau=0.5/F_s$ & exact analytic delay still recovered \\
T06 & double slope & $S\rightarrow2S$ & $f_b$ doubles \\
T07 & negative slope & $S<0$ & beat sign changes as derived \\
T08 & conjugate order & swap mixer convention & sign flips only \\
T09 & amplitude scaling & $\alpha=0.01,1,10$ noise-free & frequency unchanged \\
T10 & initial phase & random $\phi_0$ & frequency unchanged \\
T11 & two-way no offset & $\theta=0$ & $f_{AB}=f_{BA}$ \\
T12 & two-way clock offset & $\theta=100$ ps & exact $\hat\tau,\hat\theta$ \\
T13 & negative clock offset & $\theta<0$ & sign recovered \\
T14 & zero distance, nonzero offset & $\tau=0$ & two-way algebra isolates $\theta$ \\
T15 & CFO up-chirp & $\Delta f\ne0$ & beat bias equals CFO \\
T16 & CFO up/down & two slopes & recover $\Delta f,\tau$ \\
T17 & slope mismatch & $\Delta S\ne0$ & instantaneous beat has linear drift \\
T18 & clock skew & $\epsilon\ne0$ & beat drift matches first-order model \\
T19 & AWGN zero & SNR $=\infty$ & identical to ideal baseline \\
T20 & AWGN scaling & SNR sweep & variance decreases with SNR \\
T21 & common PN & same PN TX/LO & range-correlation cancellation visible \\
T22 & independent PN & independent oscillators & larger residual than common case \\
T23 & systematic curvature zero & coefficient $=0$ & exact baseline \\
T24 & sinusoidal ramp ripple & known ripple & expected beat sidebands/residual \\
T25 & group-delay symmetry & equal paths & two-way bias common-mode \\
T26 & group-delay asymmetry & unequal TX/RX & predicted clock-offset bias \\
T27 & one multipath echo & known $\alpha,\Delta\tau$ & multi-tone/vector effect visible \\
T28 & ADC quantization disabled & infinite bits & baseline preserved \\
T29 & clipping & low full scale & estimator quality flag trips \\
T30 & code identity & two orthogonal codes & correct station selected \\
T31 & code delay & fractional chip shift & naive correlation degrades \\
T32 & code alignment compensation & same T31 & recovery improves \\
\bottomrule
\end{longtable}

\section{First demo pseudocode}
\begin{verbatim}
% demo01_single_link_ideal.m
cfg = cfg_awr_smoke_v1();
cfg.link.tau = 5e-9;
cfg.impair = disable_all_impairments(cfg.impair);

t = (0:cfg.N-1).' / cfg.Fs;
phi_tx = fmcw_phase(t, cfg.waveform);
phi_rx = fmcw_phase(t-cfg.link.tau, cfg.waveform);
tx = exp(1j*phi_tx);
rx = exp(1j*phi_rx);
beat = dechirp(tx, rx, cfg.convention);

fb_theory = cfg.S * cfg.link.tau;
res_phase = estimate_phase_slope(beat, cfg.Fs);
res_fft   = estimate_fft_peak(beat, cfg.Fs);

assert(abs(res_phase.f - fb_theory) < cfg.test.freq_tol);
plot_chirp_geometry(...);
plot_beat_spectrum(...);

disp(result_table(fb_theory,res_phase,res_fft));
\end{verbatim}
The exact function names can change, but the separation between theory, waveform, estimator, assertion, and plotting should remain.

\section{Two-way demo pseudocode}
\begin{verbatim}
% demo02_twtt_ideal.m
cfg.link.tau = 5e-9;
cfg.clock.theta = 100e-12;

dAB = cfg.link.tau + cfg.clock.theta;
dBA = cfg.link.tau - cfg.clock.theta;

zAB = simulate_ideal_link(cfg, dAB);
zBA = simulate_ideal_link(cfg, dBA);

fAB = estimate_phase_slope(zAB,cfg.Fs).f;
fBA = estimate_phase_slope(zBA,cfg.Fs).f;

[tau_hat,theta_hat] = solve_twtt(fAB,fBA,cfg.S);
assert_close(tau_hat,cfg.link.tau);
assert_close(theta_hat,cfg.clock.theta);
\end{verbatim}

\section{Estimator benchmark protocol}
Use identical synthetic beat records for every estimator. For each SNR/beat frequency:
\begin{enumerate}
\item generate a clean complex tone or dechirped waveform;
\item add one frozen random-noise realization per trial and feed the same record to all estimators;
\item log frequency bias, frequency standard deviation, timing-equivalent error, phase error, failures, and runtime;
\item compare against a white-noise CRLB using the exact SNR convention;
\item repeat for off-bin frequencies uniformly distributed within one FFT bin to avoid favoring a bin-centered case.
\end{enumerate}

Candidate estimators:
\begin{itemize}
\item nearest FFT bin (baseline only);
\item parabolic/interpolated FFT peak;
\item phase-slope least squares;
\item CZT refinement;
\item direct nonlinear complex-sinusoid fit;
\item multi-chirp joint fit later.
\end{itemize}

\section{Monte Carlo experiment matrix}
A compact initial matrix is:
\begin{longtable}{p{0.21\textwidth}p{0.31\textwidth}p{0.38\textwidth}}
\toprule Sweep & Values & Output \\ \midrule
SNR & -10:5:40 dB & bias/std/RMSE for $\tau,\theta$ \\
clock offset & 0, 10, 100 ps, 1 ns & linearity and sign \\
slope & 10, 30, 60 MHz/us & timing sensitivity vs IF frequency \\
observation $N$ & 128,256,512,1024 & estimator variance vs coherent time \\
CFO & 0, 1, 10, 100 kHz & bias for single slope; recovery with up/down \\
slope mismatch & 0, 1, 10, 100 ppm & tone-to-chirp transition and bias \\
PN level & nominal mask +/- 5/10 dB & timing precision and residual PSD \\
group asymmetry & 0, 10, 100 ps & absolute TWTT bias \\
multipath & -10 to -40 dB, variable excess delay & phase/frequency bias \\
\bottomrule
\end{longtable}
Do not run the entire cross-product at first; use one-dimensional sweeps with all other impairments disabled, then selected interactions.

\section{AWR2944 config extraction checklist}
When importing the current radar package profile, store both raw config values and derived SI values:
\begin{itemize}
\item raw start-frequency register/config value and derived Hz;
\item raw slope value and derived Hz/s;
\item ADC sample rate in ksps and Hz;
\item ADC start time, ramp end time, idle time;
\item samples/chirp, chirps/frame, frames, frame period;
\item RX/TX masks and TX sequence;
\item real/complex ADC format and I/Q convention;
\item HPF corner selections;
\item calibration/monitor settings that may affect repeatability;
\item reference clock source and frequency;
\item any trigger/sync mode used.
\end{itemize}
Write these into the saved result metadata so a plot always states exactly which hardware profile it represents.

\section{DCA1000/raw-data considerations}
The DCA1000 documentation reminds us that raw ADC data is captured over UDP and then reordered/zero-filled by tooling. For a precision timing experiment, packet loss cannot be treated as benign if zeros land inside the chirp used for phase estimation. The hardware pipeline should therefore record packet-loss/reorder statistics and reject or flag corrupted chirps rather than silently averaging them.

The parser must also preserve the exact I/Q ordering and RX-channel mapping. A swapped I/Q convention conjugates/sign-flips phase and can masquerade as a timing-sign bug.

\section{Low-frequency IF feasibility checklist}
Before choosing a 10-ps-centered beat near DC on hardware, verify:
\begin{enumerate}
\item configured analog HPF corner frequency;
\item digital DC/static-clutter removal in any firmware path;
\item RX DC offset/leakage magnitude;
\item usable IF passband and ADC window;
\item whether a known intentional frequency offset can move the timing beat to a cleaner band;
\item stability/knowledge of that offset across the two units.
\end{enumerate}
A simulation option $f_{off}$ should exist from the beginning even if set to zero in V0.

\section{Two-board hardware characterization plan}
Before attempting time transfer, run a characterization ladder:
\subsection*{H0 - independent monostatic sanity}
Each board separately measures the same static corner reflector. Record range spectrum, beat phase across chirps, and temperature/power-cycle repeatability.

\subsection*{H1 - reference-clock experiment}
If feasible from TI documentation/hardware, drive both with a common external reference. Measure whether chirp beat/frequency statistics improve compared with independent references. Do not assume carrier phase is coherent.

\subsection*{H2 - trigger alignment experiment}
Use the available hardware/software trigger mechanism and measure chirp-start repeatability. If there is no direct time marker, infer relative start from a strong direct cross-link or external instrument.

\subsection*{H3 - one-direction cross-link}
Transmit from A and receive/dechirp at B if the radar mode permits. Measure CFO/slope mismatch and compare with V2 simulation.

\subsection*{H4 - reciprocal exchange}
Acquire both directions under the most simultaneous/controlled scheme available. Save raw complex data from both. Apply the ideal two-way solve and quantify bias/std.

\subsection*{H5 - calibration/repeatability}
Reverse cables/positions where meaningful, power cycle, vary temperature/time, and determine which fixed delays are stable enough to calibrate.

\section{Result-file schema}
Save a MATLAB struct roughly like:
\begin{verbatim}
result.meta.timestamp
result.meta.git_commit
result.meta.matlab_version
result.cfg.raw_hardware
result.cfg.derived_si
result.truth.tau
result.truth.theta
result.signal.beat_AB
result.signal.beat_BA
result.estimator.name
result.estimate.fAB
result.estimate.fBA
result.estimate.tau
result.estimate.theta
result.quality.snr_db
result.quality.phase_fit_rms
result.quality.packet_loss
result.random.seed
\end{verbatim}
For large Monte Carlo data, save aggregate statistics plus seeds/config rather than every waveform unless needed.

\section{Plot style for Saeed/program manager}
The first figures should answer one question each and have the theory printed on the figure or caption. Recommended visual hierarchy:
\begin{enumerate}
\item \textbf{Concept plot:} chirp slope + delay -> frequency separation.
\item \textbf{Validation plot:} measured/simulated beat peak versus theoretical $S\tau$.
\item \textbf{TWTT plot:} two beats -> recovered TOF and clock offset.
\end{enumerate}
Avoid range-Doppler heatmaps or MIMO angle plots unless they directly support the timing story. Those are familiar radar visuals but distract from the novelty direction.

\section{Software-quality gates}
\begin{itemize}
\item Every demo script starts from a clean workspace and deterministic config.
\item Every impairment has an enable flag and zero-effect test.
\item No magic unit conversions inside equations.
\item All random generators use recorded seeds.
\item Plotting functions never alter data/estimates.
\item All estimators return quality metrics, not only a scalar.
\item A regression test stores several known configurations and expected estimates.
\item Functions avoid dependence on unavailable toolboxes unless explicitly documented.
\end{itemize}

\section{Milestone schedule from here}
\begin{description}
\item[M0 - corpus complete.] Done with this archive: source map, equations, blueprint, deduped PDFs, and missing-source status.
\item[M1 - ideal single-link.] MATLAB V0 + unit tests + two clean plots.
\item[M2 - ideal TWTT.] reciprocal V1 + recovered $\tau,\theta$ plot for Saeed.
\item[M3 - independent clocks.] CFO/up-down + skew/slope mismatch.
\item[M4 - precision baseline.] AWGN estimator benchmark and CRLB comparison.
\item[M5 - physical realism.] phase noise, nonlinearity, IF/group delay, ADC.
\item[M6 - AWR mapping.] config-driven simulation matched to actual capture settings.
\item[M7 - two-board data.] calibrate/fit model against hardware.
\item[M8 - coded extension.] node identity / phase-coded FMCW only after M7 backbone is stable.
\end{description}

\section{Stop conditions and debugging order}
If an estimate is wrong, debug in this order:
\begin{enumerate}
\item units ($\mu$s versus s, MHz/us versus Hz/s);
\item mixer sign/conjugation;
\item delay convention (one-way versus round trip);
\item sample/time vector origin and ADC start time;
\item estimator phase unwrap/fit;
\item CFO/slope terms;
\item only then stochastic noise/nonlinearity.
\end{enumerate}
Most spectacular FMCW simulation errors are bookkeeping errors before they are radar-physics errors.

\end{document}
